%% Cover letter for IEEE Transactions on Information Forensics and Security.
%% Uploaded with the manuscript in the journal's submission system.

\documentclass[12pt]{article}
\usepackage[margin=2.5cm]{geometry}
\usepackage[colorlinks=true,urlcolor=blue,linkcolor=blue]{hyperref}
\pagestyle{empty}
\setlength{\parskip}{0.6em}
\setlength{\parindent}{0pt}

\begin{document}

\noindent The Editor-in-Chief\\
\emph{IEEE Transactions on Information Forensics and Security}

\bigskip

\noindent Dear Editor,

We submit our manuscript, \emph{Knowledge-Base Poisoning of Retrieval-Augmented Generation
Under a Compromised Authorized Publisher}, for consideration as a Regular Paper.

\textbf{Fit with the journal.}
The paper is an information-security study of document authentication and provenance in a
deployed class of systems. Its subject is the trust assumption behind authenticated ingestion
for retrieval-augmented generation (RAG): what a verified signature does and does not
establish once an authorized signer is compromised, and how an answer-aggregation rule should
count parties so that one key holder cannot fabricate corroboration. The threat model,
the binding analysis of the grouping attribute, and the evaluation under an adversary that
holds a valid key are security contributions and do not depend on a particular language model.
We believe the journal's readership is the appropriate audience.

\textbf{Contribution.}
(1) We add a compromised authorized publisher to the RAG ingestion adversary model and show
that the protection measured for authenticated ingestion does not survive it:
attack-attributable success rises from 0.00 to 0.82 on SQuAD and from 0.01 to 0.83 on HotpotQA
with nothing changed but possession of one publisher key.
(2) We propose a per-principal contribution cap for isolate-then-aggregate defenses, keyed on
the signer identity recorded by the ingestion service, with a stated bound on the weight one
principal can contribute. It lowers attributable success from 0.83 to 0.59 on SQuAD at five
poisoned documents per target with clean accuracy unchanged.
(3) We attack the binding of the principal directly and show that including any
uploader-chosen attribute, even one covered by the signature, returns the rule to identifier
counting.
(4) We measure the limits: no benefit against a single poisoned document, a narrowing benefit
at high budgets, no dependence on publisher diversity, and an abstention cost.

\textbf{Limits we state in the paper.}
The main experiments use a deterministic generator and small corpora. On two locally served
open-weight models the cap moves in the same direction with no regression, but the effect is
small and we report it as unconfirmed. Prior defenses were not reimplemented, so comparisons
with them are qualitative.

\textbf{Reproducibility.}
Code, configurations, seeds, model and dataset revisions, and hashed per-query outcomes are
public, and every table, figure, and reported number is regenerated from that bundle.

\textbf{Declarations.}
The manuscript is original, has not been published, and is not under consideration elsewhere.
Earlier versions were returned without external review by MDPI \emph{Information} and by
\emph{Expert Systems with Applications}, the latter as outside its aims and scope. The
accompanying statement on previous submissions quotes both decisions and explains why the
manuscript is in scope here. All authors have approved this submission. The
authors declare no competing interests, and the research received no specific funding. Use of
an AI assistant is disclosed in the Acknowledgment section as IEEE policy requires.

\bigskip

\noindent Yours sincerely,

\medskip

\noindent Md Mazharul Islam, on behalf of all authors\\
Department of Electrical and Computer Engineering, North South University\\
Dhaka 1229, Bangladesh\\
\href{mailto:mazharul.islam1@northsouth.edu}{mazharul.islam1@northsouth.edu}

\end{document}
